\documentclass[a4paper,12pt]{article}

\usepackage[T2A]{fontenc}
\usepackage[utf8]{inputenc}
\usepackage[russian]{babel}

\usepackage{amsmath, amssymb, amsthm, mathtools}
\usepackage{helvet}
\renewcommand{\familydefault}{\sfdefault}
\usepackage{icomma}
\usepackage{geometry}
\usepackage{microtype}
\usepackage{tikz}
\usepackage{tkz-euclide}
\usepackage{pgfplots}
\pgfplotsset{compat=1.18}
\usepackage{tabularx, booktabs, longtable}
\usepackage{enumitem}
\usepackage{hyperref}
\usepackage{parskip}
\usepackage{fancyhdr}
\usepackage{setspace}
\usepackage{xcolor}

\geometry{
  a4paper,
  left=18mm,
  right=18mm,
  top=18mm,
  bottom=20mm
}

\definecolor{accent}{HTML}{008080}
\definecolor{darkaccent}{HTML}{004D4D}
\definecolor{textgray}{HTML}{333333}
\definecolor{softgray}{HTML}{F2F5F5}

\hypersetup{
  colorlinks=true,
  linkcolor=darkaccent,
  urlcolor=darkaccent
}

\onehalfspacing
\setlength{\parindent}{0pt}
\setlist[itemize]{leftmargin=1.2em, itemsep=0.2em, topsep=0.3em}
\setlist[enumerate]{leftmargin=1.4em, itemsep=0.3em, topsep=0.3em}

\pagestyle{fancy}
\fancyhf{}
\lhead{\textcolor{textgray}{Чек-лист: неравенства}}
\rhead{\textcolor{textgray}{ЕГЭ}}
\cfoot{\textcolor{textgray}{\thepage}}
\renewcommand{\headrulewidth}{0.4pt}
\renewcommand{\headrule}{\hbox to\headwidth{\color{accent}\leaders\hrule height \headrulewidth\hfill}}

\newcommand{\sectiontitle}[1]{%
  \vspace{0.6em}
  {\Large\bfseries\textcolor{darkaccent}{#1}}\par
  \vspace{0.25em}
  {\color{accent}\rule{\linewidth}{0.5pt}}
  \vspace{0.4em}
}

\newcommand{\subsectiontitle}[1]{%
  \vspace{0.5em}
  {\large\bfseries\textcolor{darkaccent}{#1}}\par
  \vspace{0.15em}
}

\newcommand{\important}[1]{\textbf{\textcolor{darkaccent}{#1}}}

\newcommand{\openpoint}[2]{\draw[accent, line width=0.8pt, fill=white] (#1,0) circle (2.2pt) node[below=6pt, text=black] {$#2$};}
\newcommand{\closedpoint}[2]{\filldraw[accent] (#1,0) circle (2.2pt) node[below=6pt, text=black] {$#2$};}

\newcommand{\axisbase}[2]{%
\begin{tikzpicture}[scale=0.82, baseline=-0.5ex]
  \draw[->, line width=0.6pt] (#1,0)--(#2,0) node[right] {$x$};
  \foreach \x in {#1,...,#2} {
    \draw[black!35] (\x,0.06)--(\x,-0.06);
  }
}

\newcommand{\axisend}{\end{tikzpicture}}

\begin{document}

\begin{titlepage}
\thispagestyle{empty}
\begin{center}
\vspace*{28mm}

{\Huge\bfseries\textcolor{darkaccent}{Чек-лист}}\par
\vspace{6mm}
{\LARGE\bfseries Неравенства: запись промежутков, линейные неравенства и метод интервалов}\par

\vspace{18mm}

{\Large Лёвин Артём Александрович}\par
\vspace{2mm}
{\Large эксклюзивно для Тимофея Васильченко}\par

\vspace{12mm}

{\large ЕГЭ}\par
\vspace{4mm}
{\large \today}

\vfill

\begin{tikzpicture}[remember picture, overlay]
  \draw[accent, line width=1.1pt]
    ([xshift=18mm,yshift=24mm]current page.south west)
    .. controls ([xshift=55mm,yshift=48mm]current page.south west)
    and ([xshift=-65mm,yshift=8mm]current page.south east)
    .. ([xshift=-18mm,yshift=31mm]current page.south east);
  \draw[darkaccent, line width=0.6pt]
    ([xshift=30mm,yshift=17mm]current page.south west)
    .. controls ([xshift=72mm,yshift=35mm]current page.south west)
    and ([xshift=-92mm,yshift=46mm]current page.south east)
    .. ([xshift=-28mm,yshift=18mm]current page.south east);
\end{tikzpicture}

\end{center}
\end{titlepage}

\sectiontitle{1. Азбука неравенств и промежутков}

\subsectiontitle{1.1. Как читать простые неравенства}

\begin{center}
\renewcommand{\arraystretch}{1.35}
\begin{tabularx}{\linewidth}{@{}>{\bfseries}c X X@{}}
\toprule
Запись & Как читать & Промежуток \\
\midrule
$x\ge 5$ & $x$ не меньше $5$; подходят $5$ и все числа больше $5$ & $[5;+\infty)$ \\
$x>5$ & $x$ больше $5$; сама точка $5$ не подходит & $(5;+\infty)$ \\
$x\le 5$ & $x$ не больше $5$; подходят $5$ и все числа меньше $5$ & $(-\infty;5]$ \\
$x<5$ & $x$ меньше $5$; сама точка $5$ не подходит & $(-\infty;5)$ \\
\bottomrule
\end{tabularx}
\end{center}

\important{Главная идея.}
Если неравенство нестрогое, то граничная точка входит в ответ: точка закрашена, скобка квадратная.  
Если неравенство строгое, то граничная точка не входит в ответ: точка выколота, скобка круглая.

\subsectiontitle{1.2. Как это выглядит на числовой прямой}

\begin{center}
\begin{tabularx}{\linewidth}{@{}c X@{}}
\toprule
Неравенство & Числовая прямая \\
\midrule
$x\ge 5$ &
\axisbase{2}{8}
  \draw[accent, line width=1.8pt, ->] (5,0)--(8,0);
  \closedpoint{5}{5}
\axisend
\\[1.4em]
$x>5$ &
\axisbase{2}{8}
  \draw[accent, line width=1.8pt, ->] (5,0)--(8,0);
  \openpoint{5}{5}
\axisend
\\[1.4em]
$x\le 5$ &
\axisbase{2}{8}
  \draw[accent, line width=1.8pt, <-] (2,0)--(5,0);
  \closedpoint{5}{5}
\axisend
\\[1.4em]
$x<5$ &
\axisbase{2}{8}
  \draw[accent, line width=1.8pt, <-] (2,0)--(5,0);
  \openpoint{5}{5}
\axisend
\\
\bottomrule
\end{tabularx}
\end{center}

\important{Обязательная привычка.}
Ось нужно подписывать: $x$, $t$, $a$ или другая переменная. В задачах повышенной сложности могут встречаться разные оси.

\sectiontitle{2. Линейные неравенства}

\subsectiontitle{2.1. Что называется линейным неравенством}

Линейное неравенство — это неравенство, которое после упрощения содержит переменную только в первой степени:
\[
ax+b>0,\qquad ax+b\ge 0,\qquad ax+b<0,\qquad ax+b\le 0.
\]

Здесь нет квадратов, корней, логарифмов, дробей с переменной в знаменателе и других усложнений.

\subsectiontitle{2.2. Главное правило про отрицательное число}

\[
\boxed{\text{При умножении или делении неравенства на отрицательное число знак меняется.}}
\]

Например:
\[
-2x-1\ge 0.
\]
Переносим $-1$ вправо:
\[
-2x\ge 1.
\]
Делим на $-2$, поэтому знак меняется:
\[
x\le -\frac12.
\]

\important{Типичная ошибка.}
Ученик делит на отрицательное число, но оставляет знак неравенства прежним. Это сразу меняет ответ на противоположный.

\subsectiontitle{2.3. Алгоритм решения линейного неравенства}

\begin{enumerate}
  \item Раскройте скобки, если они есть.
  \item Перенесите все слагаемые с переменной в одну сторону, числа — в другую.
  \item Приведите подобные.
  \item Уедините переменную.
  \item Если делите или умножаете на отрицательное число, поменяйте знак неравенства.
  \item Запишите ответ промежутком и, если нужно, изобразите его на оси.
\end{enumerate}

\subsectiontitle{2.4. Пример}

Решим неравенство:
\[
-2x+3\ge 4(x+5).
\]
Раскрываем скобки:
\[
-2x+3\ge 4x+20.
\]
Переносим всё влево:
\[
-2x-4x+3-20\ge 0.
\]
Получаем:
\[
-6x-17\ge 0.
\]
Тогда
\[
-6x\ge 17.
\]
Делим на $-6$ и меняем знак:
\[
x\le -\frac{17}{6}.
\]
Ответ:
\[
x\in\left(-\infty;-\frac{17}{6}\right].
\]

\sectiontitle{3. Нормализованный вид неравенства}

Перед применением метода интервалов неравенство желательно привести к виду:
\[
F(x)\ge 0,\qquad F(x)>0,\qquad F(x)\le 0,\qquad F(x)<0.
\]

\important{Нормализованный вид} — это форма, где справа стоит $0$, а слева находится выражение, с которым дальше удобно работать.

\subsectiontitle{Важно}

\begin{itemize}
  \item Подстановку для определения знака делают именно в левую часть нормализованного неравенства.
  \item Нельзя писать конструкции вида
  \[
  F(x)<0=0.
  \]
  Это математически бессмысленная запись.
  \item После нормализации неравенство временно заменяют уравнением $F(x)=0$, чтобы найти граничные точки.
\end{itemize}

\sectiontitle{4. Метод интервалов}

\subsectiontitle{4.1. Когда нужен метод интервалов}

Метод интервалов нужен, когда неравенство уже не является линейным или приводится к произведению множителей:
\[
(x-a)(x-b)(x-c)\ge 0,
\]
а также в более сложных алгебраических неравенствах.

В этом чек-листе рассматривается базовая версия метода интервалов для многочленов и произведений линейных множителей.

\subsectiontitle{4.2. Алгоритм метода интервалов}

\begin{enumerate}
  \item Приведите неравенство к нормализованному виду: справа должен быть $0$.
  \item Замените неравенство уравнением:
  \[
  F(x)=0.
  \]
  \item Найдите корни уравнения.
  \item Отметьте корни на числовой прямой в порядке возрастания.
  \item Закрасьте точку, если неравенство нестрогое: $\ge$ или $\le$.
  \item Выколите точку, если неравенство строгое: $>$ или $<$.
  \item Возьмите удобную пробную точку из одного интервала.
  \item Подставьте её в $F(x)$ и определите знак.
  \item Расставьте знаки на остальных интервалах.
  \item Выберите интервалы, соответствующие знаку исходного неравенства.
\end{enumerate}

\subsectiontitle{4.3. Пример: квадратное неравенство}

Решим:
\[
x^2-x-6\ge 0.
\]
Заменяем неравенство уравнением:
\[
x^2-x-6=0.
\]
Раскладываем на множители:
\[
(x-3)(x+2)=0.
\]
Корни:
\[
x=-2,\qquad x=3.
\]

Так как неравенство нестрогое, точки $-2$ и $3$ закрашены.

\begin{center}
\begin{tikzpicture}[scale=0.9]
  \draw[->, line width=0.6pt] (-4.5,0)--(5.2,0) node[right] {$x$};
  \closedpoint{-2}{-2}
  \closedpoint{3}{3}
  \node[above=6pt] at (-3.3,0) {$+$};
  \node[above=6pt] at (0.5,0) {$-$};
  \node[above=6pt] at (4.2,0) {$+$};
  \draw[accent, line width=1.5pt] (-4.4,0.18)--(-2,0.18);
  \draw[accent, line width=1.5pt] (3,0.18)--(5.0,0.18);
\end{tikzpicture}
\end{center}

Нам нужны значения, при которых левая часть неотрицательна:
\[
x^2-x-6\ge 0.
\]
Значит, выбираем интервалы со знаком $+$ и сами корни:
\[
x\in(-\infty;-2]\cup[3;+\infty).
\]

\subsectiontitle{4.4. Пример: произведение четырёх множителей}

Решим:
\[
(x+10)(x-3)(x+2)(x-15)<0.
\]
Находим корни:
\[
x=-10,\qquad x=-2,\qquad x=3,\qquad x=15.
\]

Неравенство строгое, поэтому все точки выколоты.

Возьмём $x=0$:
\[
(0+10)(0-3)(0+2)(0-15)>0,
\]
потому что знаки множителей:
\[
(+)\cdot(-)\cdot(+)\cdot(-)=+.
\]

Дальше знаки чередуются:

\begin{center}
\begin{tikzpicture}[scale=0.82]
  \draw[->, line width=0.6pt] (-11.5,0)--(16.5,0) node[right] {$x$};
  \openpoint{-10}{-10}
  \openpoint{-2}{-2}
  \openpoint{3}{3}
  \openpoint{15}{15}
  \node[above=6pt] at (-10.9,0) {$+$};
  \node[above=6pt] at (-6,0) {$-$};
  \node[above=6pt] at (0.5,0) {$+$};
  \node[above=6pt] at (9,0) {$-$};
  \node[above=6pt] at (15.8,0) {$+$};
  \draw[accent, line width=1.5pt] (-10,0.18)--(-2,0.18);
  \draw[accent, line width=1.5pt] (3,0.18)--(15,0.18);
\end{tikzpicture}
\end{center}

Нужны отрицательные значения, поэтому:
\[
x\in(-10;-2)\cup(3;15).
\]

\sectiontitle{5. Если уравнение не имеет корней}

Рассмотрим неравенство:
\[
x^2-x+6>0.
\]
Связанное уравнение:
\[
x^2-x+6=0.
\]
Дискриминант:
\[
D=(-1)^2-4\cdot1\cdot6=1-24=-23<0.
\]
Корней нет.

Но это не значит, что неравенство не имеет решений. Нужно продолжить рассуждение.

Так как корней нет, знак выражения на всей числовой прямой один и тот же. Проверим, например, $x=0$:
\[
0^2-0+6=6>0.
\]
Значит, выражение положительно при всех $x$.

\[
x^2-x+6>0
\quad\Rightarrow\quad
x\in(-\infty;+\infty).
\]

Если бы было:
\[
x^2-x+6\le 0,
\]
то решений не было бы:
\[
x\in\emptyset.
\]

\important{Ключевая мысль.}
Отсутствие корней у уравнения $F(x)=0$ не означает автоматически отсутствие решений у неравенства. Нужно определить знак $F(x)$ на всей оси.

\sectiontitle{6. Кратность корня и смена знака}

\subsectiontitle{6.1. Что такое кратность корня}

Если множитель повторяется несколько раз, корень имеет кратность.

Например:
\[
(x-2)^2=(x-2)(x-2).
\]
Корень $x=2$ встречается два раза. Значит, это корень чётной кратности.

А в выражении
\[
(x+3)^3
\]
корень $x=-3$ встречается три раза. Значит, это корень нечётной кратности.

\subsectiontitle{6.2. Правило смены знака}

\[
\boxed{
\begin{array}{l}
\text{Если корень имеет нечётную кратность, знак при переходе через него меняется.}\\
\text{Если корень имеет чётную кратность, знак при переходе через него не меняется.}
\end{array}
}
\]

\subsectiontitle{6.3. Пример с чётной кратностью}

Решим:
\[
(x-2)^2(x+3)\ge 0.
\]

Корни:
\[
x=2 \quad \text{кратности }2,\qquad x=-3 \quad \text{кратности }1.
\]

Точка $-3$ имеет нечётную кратность, поэтому знак через неё меняется.  
Точка $2$ имеет чётную кратность, поэтому знак через неё не меняется.

Проверим знак на правом интервале, например при $x=3$:
\[
(3-2)^2(3+3)>0.
\]

Схема знаков:

\begin{center}
\begin{tikzpicture}[scale=0.9]
  \draw[->, line width=0.6pt] (-5,0)--(4.5,0) node[right] {$x$};
  \closedpoint{-3}{-3}
  \closedpoint{2}{2}
  \node[above=6pt] at (-4.1,0) {$-$};
  \node[above=6pt] at (-0.5,0) {$+$};
  \node[above=6pt] at (3.2,0) {$+$};
  \draw[accent, line width=1.5pt] (-3,0.18)--(4.3,0.18);
  \draw[accent, line width=0.9pt, dashed] (2,0.32) arc[start angle=180,end angle=0,x radius=0.35,y radius=0.25];
\end{tikzpicture}
\end{center}

Ответ:
\[
x\in[-3;+\infty).
\]

\subsectiontitle{6.4. Пример с большой степенью}

Рассмотрим:
\[
(x-2)^{1000}(x+3)^{1001}\le 0.
\]

Корни:
\[
x=2 \quad \text{кратности }1000,
\qquad
x=-3 \quad \text{кратности }1001.
\]

\begin{itemize}
  \item $1000$ — чётное число, значит, при переходе через $2$ знак не меняется.
  \item $1001$ — нечётное число, значит, при переходе через $-3$ знак меняется.
\end{itemize}

Проверим знак справа от всех корней, например при $x=4$:
\[
(4-2)^{1000}(4+3)^{1001}>0.
\]

Тогда схема знаков:

\begin{center}
\begin{tikzpicture}[scale=0.9]
  \draw[->, line width=0.6pt] (-5,0)--(4.5,0) node[right] {$x$};
  \closedpoint{-3}{-3}
  \closedpoint{2}{2}
  \node[above=6pt] at (-4.1,0) {$-$};
  \node[above=6pt] at (-0.5,0) {$+$};
  \node[above=6pt] at (3.2,0) {$+$};
  \draw[accent, line width=1.5pt] (-5,0.18)--(-3,0.18);
  \filldraw[accent] (2,0) circle (2.2pt);
\end{tikzpicture}
\end{center}

Так как нужно $\le 0$, берём отрицательный интервал и точки, где выражение равно нулю:
\[
x\in(-\infty;-3]\cup\{2\}.
\]

\important{Почему $2$ записана отдельно?}
Возле $2$ знак положительный с обеих сторон, но в самой точке выражение равно нулю. Поэтому $2$ подходит только как отдельная точка.

\sectiontitle{7. Типичные ошибки}

\begin{enumerate}
  \item \important{Не подписывать ось.} В задачах могут быть оси $x$, $t$, $a$ и другие.
  \item \important{Путать строгие и нестрогие неравенства.} Знак $>$ или $<$ даёт выколотую точку; знак $\ge$ или $\le$ даёт закрашенную.
  \item \important{Забывать менять знак при делении на отрицательное число.}
  \item \important{Подставлять пробную точку не в то выражение.} Подстановка делается в левую часть нормализованного неравенства.
  \item \important{Думать, что если уравнение не имеет корней, то неравенство тоже не имеет решений.} Это неверно.
  \item \important{Всегда чередовать знаки без учёта кратности.} При чётной кратности знак не меняется.
  \item \important{Брать $0$ как пробную точку, если $0$ является корнем.} В этом случае нужно взять другое удобное число.
  \item \important{Неверно записывать одиночную точку.} Отдельная точка записывается через фигурные скобки, например $\{2\}$.
\end{enumerate}

\sectiontitle{8. Мини-чек-лист перед ответом}

Перед тем как записать ответ, проверьте:

\begin{itemize}
  \item справа действительно стоит $0$;
  \item все корни найдены;
  \item корни расположены на оси в порядке возрастания;
  \item точки закрашены или выколоты в соответствии со знаком неравенства;
  \item кратности корней учтены;
  \item знак выбран именно тот, который требует неравенство;
  \item ответ записан промежутками корректно;
  \item бесконечности записаны только с круглыми скобками.
\end{itemize}

\newpage

\sectiontitle{Тренировочный блок}

Решите неравенства. Ответы запишите промежутками.

\subsectiontitle{Средний уровень}

\begin{enumerate}
  \item Решите неравенство:
  \[
  7x-14\ge 0.
  \]

  \item Решите неравенство:
  \[
  -5x+20<0.
  \]

  \item Решите неравенство:
  \[
  3(x-2)\le 2x+5.
  \]
\end{enumerate}

\subsectiontitle{Выше среднего уровня}

\begin{enumerate}[resume]
  \item Решите неравенство:
  \[
  x^2-5x+6\ge 0.
  \]

  \item Решите неравенство:
  \[
  -x^2+3x+4>0.
  \]

  \item Решите неравенство:
  \[
  (x+4)(x-1)(x-6)<0.
  \]

  \item Решите неравенство:
  \[
  (x-3)^2(x+1)\ge 0.
  \]

  \item Решите неравенство:
  \[
  x^2+2x+5\le 0.
  \]

  \item Решите неравенство:
  \[
  (x+2)^4(x-5)^3>0.
  \]
\end{enumerate}

\subsectiontitle{Сложный уровень}

\begin{enumerate}[resume]
  \item Решите неравенство:
  \[
  (x+3)^2(x-1)^5(x-4)^4\le 0.
  \]
\end{enumerate}

\newpage

\sectiontitle{Ответы к тренировочному блоку}

\begin{table}[h!]
\centering
\caption{Ответы}
\renewcommand{\arraystretch}{1.45}
\begin{tabularx}{\linewidth}{@{}c X@{}}
\toprule
№ & Ответ \\
\midrule
1 & $x\in[2;+\infty)$ \\
2 & $x\in(4;+\infty)$ \\
3 & $x\in(-\infty;11]$ \\
4 & $x\in(-\infty;2]\cup[3;+\infty)$ \\
5 & $x\in(-1;4)$ \\
6 & $x\in(-\infty;-4)\cup(1;6)$ \\
7 & $x\in[-1;+\infty)$ \\
8 & $x\in\emptyset$ \\
9 & $x\in(5;+\infty)$ \\
10 & $x\in[-3;1]$ \\
\bottomrule
\end{tabularx}
\end{table}

\end{document}